\section{Coverage estimation on the nested grid}
\label{sec:application-validation}

\subsection{Spatial support of the coverage measures}

This extension was specified on 7 September 2026 before calculating its
coverage contrasts or comparison errors. It uses 3,036 synchronous daily
fields in the 33 evaluation summers and retains the primary-grid peak time
and original event labels. The dense fields contain 465 sites, including all
121 original sites and 344 additional locations. Embedded values agree with
the primary fields to within $10^{-6}$ degrees Celsius. Every summer has
92 complete daily fields. The original graph contrast is reproduced with
absolute numerical error below $10^{-10}$.

Each site represents a cell extending $0.4^\circ$ in longitude and
$0.45^\circ$ in latitude from its centre. Cells are clipped to
$[105,125]^\circ$E by $[20,42]^\circ$N and the union of land polygons in the
Natural Earth low-resolution map distributed with R package \texttt{spData}.
All countries within the original rectangle are included. Retained cells
represent 3,220,939 km$^2$, or 95.69\% of the 3,366,047 km$^2$ of mapped land
in the rectangle. The denominator uses the represented land. All 465 cells
have positive land area. Area is computed on the WGS84 ellipsoid, including
polygon holes. A coastal cell has one observed temperature and its clipped
land area. Grid adjacency gives connectivity at this resolution; subcell
coastline connectivity is unresolved.

Four-neighbour components use adjacent grid rows or columns; eight-neighbour
components additionally join diagonal neighbours. Missing lattice positions
remain absent. For each day and threshold, both measured curves satisfy
$0\le C_t(u)\le A_t(u)\le1$ and decrease monotonically with $u$. The
additional-site exceedance proportion uses the 344 new cells and their own
normalised land areas. Connected area uses the complete retained lattice.

\subsection{Threshold dependence of coverage}

The pooled 465-site values from 2015 and 2022 determine the temperature grid.
Their 1st and 99.5th sample percentiles are rounded outward to multiples of
0.25$^\circ$C, producing 9.5--28.5$^\circ$C. All 77 thresholds use strict
exceedance, $y_{it}>u$. These levels describe the temperature field.
For $F=A$ or $C$, the contrast is
\[
 \Delta_F(u)=\frac{1}{33}\sum_{y\in\mathcal Y_H}
 \frac{1}{3}\sum_{m\in\{6,7,8\}}
 \left\{\frac{\sum_{t\in\mathcal E_{ym}}F_t(u)}{|\mathcal E_{ym}|}
 -\frac{\sum_{t\in\mathcal M_{ym}}F_t(u)}{|\mathcal M_{ym}|}\right\}.
\]
All annual curves and both group means are retained. Square-kilometre
contrasts multiply these differences by the represented land area.

Descriptive uncertainty uses 1,999 circular moving-block resamples of the
35 calendar positions from 1991 to 2025, with block length five and seed
20260907. Seven uniformly drawn block starts produce 35 positions per
resample. Development-year positions are then omitted and the remaining
annual contrasts averaged. All thresholds and outcomes share the same draws.
The 2.5th and 97.5th percentiles form pointwise intervals. Paired model-error
intervals use the same resampling and condition on the fitted
cross-validation predictions. They describe variation in the observed
sequence; they are pointwise rather than simultaneous bands.

Eight-neighbour adjacency changes the high-minus-middle connected-area
contrast by at most 0.965 percentage points over the threshold grid.
Replacing clipped land areas by normalised cosine-latitude weights changes
the total-area and connected-area contrasts by at most 0.484 and 0.611
points, respectively. The two alternating 0.5$^\circ$C subgrids retain the
overall model ordering for all-day errors.

\subsection{Coverage at the illustrative threshold}

High-regional-mean days have greater simultaneous exceedance area and larger
connected exceedance regions throughout the evaluated temperature range
(main-text Figure~4). At 24$^\circ$C, the mean exceedance
proportion rises from 37.39\% on middle days to 50.29\% on high days. The
12.91-percentage-point difference corresponds to approximately
416,000 km$^2$ (pointwise 95\% block-resampling interval:
377,000--459,000 km$^2$). The largest connected region rises from 32.10\%
to 44.67\% of represented land, a difference of 12.57 points or approximately
405,000 km$^2$ (365,000--452,000 km$^2$). Both differences are positive in all
33 summers at this threshold.


\subsection{Regression specification}

The baseline inputs are the original regional mean and day of summer. The
strong simple model adds $V_t$, the signed ordinary least-squares slope of
WBT against latitude, and the 121-site direct exceedance proportion. Each
dense cell's area is assigned to its nearest original site in the original
projected kilometre coordinates to obtain fixed primary representative
weights. Thus the direct proportion covers the same represented support
using 121 temperatures.

The graph model adds $\log\{Q_{h_k}(\mathbf y_t)/V_t\}$ for all five
original physical bandwidths. The variogram model adds the corresponding
log ratios of five binned semivariances to $V_t$. Bins use the quintiles of
the 7,260 original pair distances: internal boundaries are 600.71, 850.88,
1,201.42 and 1,553.64 km. Ties stay in the same bin; bin sizes are 1,422,
1,471, 1,454, 1,458 and 1,455. Every pair belongs to one bin. All observed
variances and energies are positive. The specified constant-field convention
sets the five log-shape inputs to zero when $V_t=0$.

Each baseline scalar has a quadratic B-spline basis with four equally spaced
knots over its training range, linear extrapolation and one redundant basis
column removed. A constant training variable is omitted from the spline basis.
Graph and variogram shapes enter linearly. All features are standardised
using training observations. Weighted ridge least squares fits an
unpenalised intercept and penalised slopes; training weights give equal
weight to months within years. Predicted proportions are clipped to $[0,1]$.
The same procedure is used for total area, connected area and the
additional-site area. One penalty per method, endpoint and outer fold is
selected from $\{0.1,1,10,100,1000\}$ by inner mean absolute error averaged
equally across thresholds, months and years.

The interpolation comparison uses a Delaunay piecewise-linear interpolant
in the original projected coordinates. Eight dense locations lie outside
the convex hull and use their nearest original site's temperature. Coverage
is then calculated directly from the interpolated field. Interpolation
preserves the original-site values to numerical precision.

The regressions estimate each threshold and outcome separately. Across all
held-out day--threshold pairs, fitted connected area exceeds fitted total
area in 0.31\%, 1.44\%, 1.59\% and 1.78\% of predictions for the mean,
simple, variogram and graph models. The maximum graph-model excess is
3.30 percentage points. Adjacent-threshold increases occur in 2.23\% of
graph area predictions and 7.49\% of graph connected-area predictions.
These discrepancies describe approximation error in the regression summaries;
measured and interpolated fields satisfy the coverage identities.

\subsection{Temporal validation}

The outer blocks are 1991--1995, 1996--2000, 2001--2005, 2006--2010,
2011--2015, 2016--2020 and 2021--2025. Each is held out in turn. Training
excludes the held-out block, its adjacent calendar years and the two
development summers. Inner validation uses the remaining calendar blocks
within outer training, again excluding adjacent years. Spline knots,
standardisation and penalties are fitted within these training sets. All
methods share the same memberships. Training may use years before and after
the held-out block; this evaluates transport across omitted years.

For method $j$, summer $y$ and outcome $F$, define
\[
 L^F_{j,y}=\frac{1}{3}\sum_{m\in\{6,7,8\}}
 \frac{1}{|\mathcal T_{ym}|}\sum_{t\in\mathcal T_{ym}}
 \frac{1}{77}\sum_{u\in\mathcal U}
 |F_t(u)-\widehat F^{(-b(y))}_{j,t}(u)|.
\]
High-day errors replace $\mathcal T_{ym}$ by $\mathcal E_{ym}$.
Paired differences subtract graph error from the strong-simple or variogram
error within each summer. Every method supplies 3,036 held-out daily curves.
Positive differences indicate lower graph error.

\begin{table}[!htbp]
\centering\small
\caption{Paired reduction in held-out mean absolute error after adding graph
information, in percentage points. Positive values favour the graph model.
Intervals use five-year calendar-block resampling of the annual paired errors,
conditional on the fitted cross-validation predictions.}
\label{tab:application-paired}
\begin{tabular}{@{}lllrrr@{}}
\toprule
Outcome & Days & Reference & Reduction & 95\% interval & Improved summers\\
\midrule
Area & all & simple & $0.0069$ & $[0.0057,0.0081]$ & 28/33 \\
Area & all & variogram & $-0.0053$ & $[-0.0070,-0.0038]$ & 4/33 \\
Area & high & simple & $0.0041$ & $[0.0005,0.0082]$ & 21/33 \\
Area & high & variogram & $-0.0017$ & $[-0.0046,0.0010]$ & 12/33 \\
Additional sites & all & simple & $0.0077$ & $[0.0062,0.0090]$ & 28/33 \\
Additional sites & all & variogram & $-0.0070$ & $[-0.0087,-0.0054]$ & 4/33 \\
Additional sites & high & simple & $0.0043$ & $[0.0006,0.0081]$ & 20/33 \\
Additional sites & high & variogram & $-0.0029$ & $[-0.0062,0.0005]$ & 12/33 \\
Connected area & all & simple & $0.0100$ & $[0.0069,0.0129]$ & 26/33 \\
Connected area & all & variogram & $-0.0137$ & $[-0.0172,-0.0101]$ & 2/33 \\
Connected area & high & simple & $0.0036$ & $[-0.0033,0.0100]$ & 18/33 \\
Connected area & high & variogram & $0.0016$ & $[-0.0028,0.0060]$ & 18/33 \\
\bottomrule
\end{tabular}
\end{table}


\begin{figure}[!t]
\centering
\includegraphics[width=\textwidth]{application_increment.pdf}
\caption{Threshold-specific paired reduction in mean absolute error from
adding graph information. Positive values favour the graph model. Blue solid
curves compare the strong simple baseline with the graph model; orange dashed
curves compare the variogram model with the graph model. Left panels use all
days and right panels use high days. Shading gives pointwise 95\% five-year
calendar-block resampling intervals, conditional on fitted predictions.
All 77 thresholds are retained.}
\label{fig:application-increment-supp}
\end{figure}
